\section{Method}
\label{sec:method}

Our overarching data curation pipeline is visualized in \cref{fig:pipeline}.
It maps Non-Euclidean raw CAD intent---comprising a structured Ops JSON and multi-view
renders---into a highly structured Euclidean execution space represented by CadQuery
Python code.

\begin{figure*}[t]
  \centering
  \includegraphics[width=\textwidth]{figure_1}
  \caption{%
    \textbf{The CADLoop Data Curation Pipeline.}
    Raw 3D geometries and semantic JSON intents are initially mapped into executable
    CadQuery code via a VLM agent cascade.
    The outputs are subsequently refined by a deterministic geometric fallback loop to
    produce high-fidelity, verified training pairs.
  }
  \label{fig:pipeline}
\end{figure*}

\subsection{Problem Formulation \& Verification Metric}
\label{sec:formulation}

The input to our curation pipeline consists of an initially generated candidate CadQuery
script $C_{init}$ (produced by an upstream VLM), a target ground-truth solid $S_{gt}$ (in
STEP format), and a reference Ops JSON $J$ providing the discrete semantic constraints
(\eg, \texttt{cut\_hole}, \texttt{tool\_size}, \texttt{centers}).
The candidate script $C_{init}$ compiles into a generated 3D boundary representation
solid $S_{gen}$.

To rigorously ensure engineering-grade accuracy, correctness is not evaluated via
heuristic distance approximations, but rather verified using exact 3D volumetric
Intersection over Union (IoU) via precise B-Rep Boolean operations computed by the
OpenCASCADE kernel:
\begin{equation}
  \text{IoU}(S_{gen}, S_{gt}) =
    \frac{\text{Vol}(S_{gen} \cap S_{gt})}
         {\text{Vol}(S_{gen}) + \text{Vol}(S_{gt}) - \text{Vol}(S_{gen} \cap S_{gt})}
  \label{eq:iou}
\end{equation}

\subsection{The Deterministic Fallback Mechanism}
\label{sec:fallback}

\Cref{fig:arch} highlights the architectural superiority of CADLoop over standard
iterative loops.
We wrap the baseline generation with a ``Geometric Constraint Discovery'' module.
When standard visual/textual validation fails, rather than blindly re-prompting the LLM,
the system triggers the targeted programmatic skill to explicitly calculate and correct
parameters.

\begin{figure*}[t]
  \centering
  \includegraphics[width=\textwidth]{figure_2}
  \caption{%
    \textbf{Architectural Comparison.}
    (A) Standard Re-Act loops rely entirely on probabilistic inference, leading to
    hallucination cycles when faced with continuous geometric errors.
    (B) CADLoop acts as an outer fallback envelope, utilizing exact geometric discovery to
    analytically modify the code's AST without VLM intervention.
  }
  \label{fig:arch}
\end{figure*}

The curation and fallback pipeline is formalized in \cref{alg:cadloop}.
By taking the generated code and reference JSON as direct inputs, we avoid a fragmented
toolset.
We expose a single unified skill (\texttt{cad\_data\_generation}) to the overall system,
which acts as an autonomous post-generation processor to analytically calculate the missing
spatial constraints and inject them directly into the Python code.

\begin{algorithm}[t]
  \caption{CADLoop Deterministic Fallback Pipeline}
  \label{alg:cadloop}
  \begin{algorithmic}[1]
    \Require Generated code $C_{init}$, reference Ops JSON $J$, GT solid $S_{gt}$,
             unified skill $\mathcal{S}_{cad}$
    \Ensure Verified CadQuery script $C^*$, or \textsc{null}
    \State $S_{gen} \leftarrow \textsc{Compile}(C_{init})$
    \State $score \leftarrow \textsc{IoU}(S_{gen},\, S_{gt})$
    \If{$score \geq 0.99$}
      \State \Return $C_{init}$ \Comment{Auto-Pass}
    \EndIf
    \If{$score < 0.80$}
      \State \Return \textsc{null} \Comment{Skip structural hallucinations}
    \EndIf
    \State \Comment{Trigger Deterministic Fallback (Proximal Near-Miss Regime)}
    \State \Comment{Skill parses $J$ to discover root cause and compute exact params}
    \State $params \leftarrow \mathcal{S}_{cad}.\textsc{CalcParams}(S_{gen},\, S_{gt},\, J)$
    \State \Comment{Deterministically update the generated code's AST}
    \State $C_{ref} \leftarrow \textsc{InjectParams}(C_{init},\, params)$
    \State $S_{ref} \leftarrow \textsc{Compile}(C_{ref})$
    \State $score' \leftarrow \textsc{IoU}(S_{ref},\, S_{gt})$
    \If{$score' \geq 0.99$}
      \State \Return $C_{ref}$
    \Else
      \State \Return $\textsc{EscalateToFlywheel}(C_{ref},\, J,\, S_{gt})$
    \EndIf
  \end{algorithmic}
\end{algorithm}

\subsection{The Unified Geometric Skill \& Parametric Grounding}
\label{sec:skill}

Rather than maintaining a scattered, hard-to-route library of narrow tools (\eg, a
specific tool for circles, another for lines), CADLoop encapsulates its deterministic
solvers into a single, cohesive programmatic skill.
Upon invocation, this skill analyzes the Ops JSON ($\mathbf{C}_{json}$) to infer the
geometric intent, dynamically instantiates the correct analytical mathematical solver, and
directly injects the explicit parameters ($\mathbf{P}_{exact}$) back into the code.
We highlight two primary internal operations this unified skill performs to resolve
failures, as visualized in \cref{fig:skill}.

\begin{figure*}[t]
  \centering
  \includegraphics[width=\textwidth]{figure_3}
  \caption{%
    \textbf{Unified Skill Internal Reasoning.}
    The tool calculates exact arc midpoints via rotation matrices and isolates inner-hole
    volume differences.
    It updates the AST deterministically, shifting the IoU score from a failing 0.893 to
    a verified 0.999.
  }
  \label{fig:skill}
\end{figure*}

\noindent\textbf{Operation 1: Arc Midpoint Resolution.}
A prevalent VLM error is specifying arcs by center and radius (\texttt{radiusArc}), which
causes severe topological inversions due to ambiguity in the sweep direction.
To resolve this, the tool parses the \texttt{Arc3D} constraints from the JSON and computes
the guaranteed unambiguous midpoint $P_{mid}$ using a 2D rotation matrix $\mathbf{R}$:
\begin{equation}
  P_{mid} = P_{center} + r \cdot \mathbf{R}(\hat{v},\, \theta/2) \cdot \vec{d}_{start}
  \label{eq:arcmid}
\end{equation}
The AST modifier then automatically replaces the brittle \texttt{radiusArc} syntax with a
robust \texttt{threePointArc} node.

\noindent\textbf{Operation 2: Feature-Level Volumetric Verification.}
To detect complex profile omissions (\eg, an LLM failing to subtract an inner bounding
loop), the tool relies on absolute volume differentials:
$\Delta V = |\text{Vol}(S_{gen}) - \text{Vol}(S_{gt})|$.
If $\Delta V$ closely matches the theoretical volume of a specific target feature from the
JSON, it triggers a localized cross-sectional area verification:
\begin{equation}
  Area_{expected} = Area_{rectangle} \pm Area_{semicircle}
  \label{eq:area}
\end{equation}
By identifying whether the area aligns with the positive or negative term of
\cref{eq:area}, the skill determines the exact structural inversion (Convex Down vs.\
Convex Up) and actively rewrites the directional Boolean constraints in the CadQuery
script.
